\subsection{Closure group, quartic bridge, and claim firewall}
\label{sec:bt1453-bt1455-group-quartic-firewall}

The closure/shear layer generated by the Szilassi involution and the D4 shear is
not dihedral \(D_9\).  On the 12-state closure space, the generators are
\[
\tau_4(b,p)=(b\oplus2,p),
\qquad
S(b,p)=(b,p+b\bmod3).
\]
The generated group has order profile
\[
1^1,\quad2^3,\quad3^8,\quad6^6,
\]
center \(C_3\), and derived subgroup \(C_3\).  Hence the closure/shear layer is
\[
S_3\times C_3.
\]

The varied golden quartic coefficient supplies a structural bridge but not a
physical derivation:
\[
4-\phi^2=3+\phi=\sqrt{13+\phi^5}.
\]
Thus the linear coefficient reads as three Szilassi opposite pairs plus a golden
tail, while its square reads as a 13-half-turn core plus a fifth-order golden
tail.  The exact finite closure remains
\[
3\cdot2\cdot2=12,\qquad 2\cdot12=24,\qquad12(13+1)=168.
\]

A claim firewall now separates exact finite facts from numerical resonances and
blocked physical claims.  Coordinate, active/guard, group, and retwined decoder
statements may be stated exactly.  Quartic and rounded-\(g\) links are resonance
claims.  Formula-level and real-world particle claims remain blocked pending
transcription and audit of Otto's rendered equations.
